\section{Problem Identification}

The literature survey highlights several limitations, a result of how these systems are designed and evaluated.
Current approaches prioritize short-term engagement and observable signals, often ignoring aspects 
like the human behavior (cognition, preferences, and satisfaction).
These limitations also lead to unintended consequences such as user fatigue, repetitive recommendations, 
and declining long-term engagement.

This section summarizes the key problems identified in the literature survey 
and motivates the need for a more human-centered recommendation framework.

%%%%%%%%%%%%%%%%%%%%%%%%%%%%%%%%%%%%%%%%%%%%%%%%%%%%%%%%%%%%%%%%
%%%%%%%%%%%%%%%%%%%%%%%%%%%%%%%%%%%%%%%%%%%%%%%%%%%%%%%%%%%%%%%%

\subsection{Short-Term Engagement Bias}

The focus on immediate user interaction has heavily influenced the design toward 
metrics like click-through rate (CTR), watch time, or the probability that a user
will interact with the next recommended item.
While these metrics are easy to measure and optimize for; they primarily represent 
short-term engagement rather than long-term satisfaction.

A recommendation that successfully triggers a click does not necessarily mean that
the user found it valuable or enjoyable; which in turn does not guarantee their 
retension over an extended period of time.
Additionally, users may interact with content out of curiosity rather than genuine interest.

Because of this, many recommendation algorithms are trained to 
maximize short-term interactions rather than long-term user value.
Over time, this can encourage content that attracts attention 
quickly but may not contribute satisfying user experiences.

This creates a bias in recommendation systems: \textbf{the system becomes optimized for immediate engagement rather than sustainable user satisfaction.}

%%%%%%%%%%%%%%%%%%%%%%%%%%%%%%%%%%%%%%%%%%%%%%%%%%%%%%%%%%%%%%%%
%%%%%%%%%%%%%%%%%%%%%%%%%%%%%%%%%%%%%%%%%%%%%%%%%%%%%%%%%%%%%%%%

\subsection{Lack of Cognitive State Modeling}

Traditional algorithms of recommendation work around 
available historical data about the user interactions.
In these models, a user is typically represented as 
a sequence of past clicks, views, or purchases.

This approach is effective for identifying behavioral patterns; however, it treats users 
as static data sources, completely ignoring their changing mental states and intentions.
Important psychological factors such as attention, curiosity, 
boredom, and fatigue does not have an explicit representation in the model.

As a result, the system can observe \textit{what users do}, but it 
cannot infer \textit{how users feel} about their current experience.
For example, a user who continues interacting with recommendations might appear 
highly engaged, even if they are becoming tired of the repetitive content.

Without modeling a cognitive state, recommender systems lack the ability to adapt 
to a strategy which can work around how users are experiencing the process.

%%%%%%%%%%%%%%%%%%%%%%%%%%%%%%%%%%%%%%%%%%%%%%%%%%%%%%%%%%%%%%%%
%%%%%%%%%%%%%%%%%%%%%%%%%%%%%%%%%%%%%%%%%%%%%%%%%%%%%%%%%%%%%%%%

\subsection{Absence of Fatigue and Diminishing Utility}

Another important limitation observed is the assumption that 
user utility remains constant across repeated interactions.
Most models assume that if a user interacted with some content in the past, 
recommending similar items repeatedly is a good move.

However, in reality, this repeatition often leads to \textit{diminishing utility}.
As users encounter similar recommendations multiple times, 
the perceived value of those recommendations decreases.
This leads to boredom, reduced interest, and eventually disengagement.

For example, if a user is recommended several videos from a single category or receives
similar product recommendations, the novelty of those recommendations gradually declines.
Despite that, many models continue to reinforce the same patterns 
because they rely heavily on past interaction signals.

Since this notion of fatigue is not explicitly modeled, recommender systems 
unintentionally generate overly repetitive recommendation streams.

%%%%%%%%%%%%%%%%%%%%%%%%%%%%%%%%%%%%%%%%%%%%%%%%%%%%%%%%%%%%%%%%
%%%%%%%%%%%%%%%%%%%%%%%%%%%%%%%%%%%%%%%%%%%%%%%%%%%%%%%%%%%%%%%%

\subsection{Static Recommendation Policies}

Most model from the traditional setting learn a single strategy representing the training data.
Therefore, given the historical pretext, these models perform really well.
Unfortunately, a user is a dynamic entity, with evolving flavours and tastes.
A single strategy is insufficient to keep track of a user's state or context.

This means the system is incapable of dynamically adjusting its decision strategy 
based on factors such as session length, user exploration patterns, or signs of fatigue.
Once a ranking model is trained, it is used to generate recommendations in the same way for every interaction.

In practice, different situations may require different recommendation strategies.
Early in a session, users may prefer familiar and relevant content.
Later in the session, introducing novelty or diversity might help prevent fatigue.

Since current systems typically rely on static ranking policies, they lack 
the ability to adapt its strategy to changing user cognitive state.

%%%%%%%%%%%%%%%%%%%%%%%%%%%%%%%%%%%%%%%%%%%%%%%%%%%%%%%%%%%%%%%%
%%%%%%%%%%%%%%%%%%%%%%%%%%%%%%%%%%%%%%%%%%%%%%%%%%%%%%%%%%%%%%%%

\subsection{Misaligned Evaluation Metrics}

Evaluation metrics play a critical role in shaping recommender system design.
Most systems are evaluated using metrics such as Precision@K, Recall@K, and 
NDCG; these metrics measure how accurately the system predicts user interactions.

These metrics are useful for benchmarking algorithm performance, however, 
only for predictive accuracy, rather than long-term user outcomes.
A system with a high ranking accuracy can still produce 
repetitive content, thus reducing user satisfaction.

Furthermore, offline evaluation methods often fail to capture the dynamic 
feedback loops that take place with real user interactions over time.
As a result, models that perform well in offline experiments 
may not always produce the best long-term user experiences.

%%%%%%%%%%%%%%%%%%%%%%%%%%%%%%%%%%%%%%%%%%%%%%%%%%%%%%%%%%%%%%%%
%%%%%%%%%%%%%%%%%%%%%%%%%%%%%%%%%%%%%%%%%%%%%%%%%%%%%%%%%%%%%%%%

\subsection{Summary of Identified Problems}

The analysis above highlights several key limitations in current recommender system design:

\begin{table}[h]
\caption{TODO Add a caption\label{tab:TODOlabel0102526}}
\centering
\begin{tabularx}{\columnwidth}{|l||Y|}
\hline
Problem & Description \\
\hline
Short-term engagement bias & Systems prioritize immediate interactions rather than long-term satisfaction \\
Lack of cognitive state modeling & Users are represented only through behavior sequences \\
Absence of fatigue modeling & Repeated recommendations may reduce user interest over time \\
Static recommendation policies & Systems do not adapt strategies based on user state \\
Misaligned evaluation metrics & Metrics reward interaction prediction rather than user wellbeing \\
\hline
\end{tabularx}
\end{table}

Together, these limitations suggest that current recommender systems are primarily optimized for \textbf{short-term behavioral signals rather than long-term human-centered outcomes}.

These challenges motivate the development of new recommendation frameworks that can explicitly consider user cognitive state, fatigue, and long-term satisfaction when generating recommendations.

